% Abstract — Route B (v3 audit arc, 2026-08-26)
\begin{abstract}
Deployment-period parameter updates for tool-using agents (agent
test-time RL) are attractive, and recent systems report gains. We show
that the most common reported gains are artifacts of three silent
failure modes, demonstrated by re-running one pipeline under
increasingly strict protocol audits: (F1) \emph{served-policy drift} ---
the updated model never becomes the serving policy, so per-task
``update effects'' are sampling-stream displacement (96/96 identical
outcomes across arms); (F2) \emph{evaluator leakage and anchor
injection} --- the hidden evaluator controls episode termination and
hand-constructed anchors seed the replay buffer, producing a phantom
$+0.070$ AUPC transfer ($p{=}0.016$) that vanishes when both are
removed; (F3) \emph{protocol-correct updates are harmful} --- under
full isolation, REINFORCE-style updates significantly degrade
first-attempt AUPC (Mistral-7B: naive $-0.19$, $p{=}0.008$, 0/8 seeds
positive; frozen deterministic across seeds). The harm replicates on a
second backbone (Qwen2.5-7B: both update arms 8/8 seeds below frozen,
$p{=}0.008$), degrades the sealed never-trained holdout, and no
alternative update rule transfers. A pre-commit gate validating
candidate adapters on accessible evidence eliminates the harm (0/8
harmful commits). As a positive control, the same serving runtime
against the official \texttt{tau2} benchmark completes both evaluated
retail tasks with full reward in 5/10 seeded runs using only a local
14B model. We release the audit chain as a reproducible harness ---
policy-consistent runtime, request-scoped RNG, hidden-evaluator
isolation, accessible-evidence utilities, integration gates --- so any
agent test-time RL pipeline can be checked for these failure modes.
\end{abstract}
